\documentclass[a4paper]{article} %可加上 twocolumn 來變成雙欄格式
\usepackage{fontspec}   %加這個就可以設定字體
\usepackage{xeCJK}       %讓中英文字體分開設置
\usepackage{indentfirst}
\usepackage{listings}
\usepackage[newfloat]{minted}
\usepackage{float}
\usepackage{graphicx}
\usepackage{caption}
\usepackage{fancyhdr}
\usepackage{hyperref}
\usepackage{amsmath}
\usepackage{multirow}
\usepackage[dvipsnames]{xcolor}
\usepackage{graphicx}
\usepackage{tabularx}
\usepackage{booktabs}
\usepackage{caption}
\usepackage{subcaption}
\usepackage{pifont}
\usepackage{amssymb}
\usepackage[backend=biber]{biblatex}
\addbibresource{main.bib}


\usepackage{pdftexcmds}
\usepackage{catchfile}
\usepackage{ifluatex}
\usepackage{ifplatform}

\usepackage[breakable, listings, skins, minted]{tcolorbox}
\usepackage{etoolbox}
\setminted{fontsize=\footnotesize}
\renewtcblisting{minted}{%
    listing engine=minted,
    minted language=python,
    listing only,
    breakable,
    enhanced,
    minted options = {
        linenos, 
        breaklines=true, 
        breakbefore=., 
        % fontsize=\footnotesize, 
        numbersep=2mm
    },
    overlay={%
        \begin{tcbclipinterior}
            \fill[gray!25] (frame.south west) rectangle ([xshift=4mm]frame.north west);
        \end{tcbclipinterior}
    }   
}

\usepackage[
top=1.5cm,
bottom=0.75cm,
left=1.5cm,
right=1.5cm,
includehead,includefoot,
heightrounded, % to avoid spurious underfull messages
]{geometry} 

\newenvironment{code}{\captionsetup{type=listing}}{}
\SetupFloatingEnvironment{listing}{name=Code}


% set the title and name
\title{Lab-1 DDPM}
\author{112550011 李佑軒}
\date{\today}


% \setCJKmainfont{Noto Serif TC}
\setCJKmainfont{Microsoft JhengHei}


\ifwindows
\setmonofont[Mapping=tex-text]{Consolas}
\fi

\XeTeXlinebreaklocale "zh"             %這兩行一定要加，中文才能自動換行
\XeTeXlinebreakskip = 0pt plus 1pt     %這兩行一定要加，中文才能自動換行

\newcommand*{\dif}{\mathop{}\!\mathrm{d}}

\begin{document}
\maketitle % write the title, name and date

\section{Introduction}


\section{Task 1 - Swiss Roll}

\subsection{Implementation details}

\subsection{q\_sample implementation} 

\subsection{Loss curve and output distribution}


\section{Task 2 - Image Generation}

\subsection{Implementation details}

\subsection{Trajectory figures and three beta schedulers}

\subsection{Result of different Predictors and discuss}

\subsection{Training and FID score}

% \begin{code}
% \captionof{listing}{\textbf{Actor Network Implementation}}
% \label{code:example}
% \begin{minted}
% class example():
%     def __init__(self):
%         pass
% \end{minted}
% \end{code}

% \begin{enumerate}
% \item \textbf{Item1}: This is the example of itemize environment. 
% \item \textbf{Item2}: This is the example of itemize environment.
% \end{enumerate}

% \subsubsection{Experimental settings}
% \begin{itemize}
% 	\item LR: 0.0001
% 	\item Optimizer: AdamW
%     \item scheduler: CosineAnnealingLR
% 	\item Batch size: 64
%     \item epochs: 64
% 	\item Seed: 77
% \end{itemize}
% \subsubsection{Experimental settings2}
% \begin{itemize}
% 	\item LR: 0.001
% 	\item Optimizer: AdamW
% 	\item scheduler: CosineAnnealingLR
% 	\item Batch size: 64
% 	\item epochs: 64
% 	\item Seed: 77
% \end{itemize}

% \begin{table}[H]
% \centering
% \caption{Example of table}
% \label{tab:efficiency_comparison}
% \begin{tabular}{lcc} % (3 columns) l: left, c: center, r: right
% \toprule
% \textbf{Algorithm} & \textbf{Score1} & \textbf{Score2} \\
% \midrule
% Algorithm1 & 16 & 30 \\
% Algorithm2 & 36 & 50 \\
% Algorithm2 & 88 & 66 \\
% \bottomrule
% \end{tabular}
% \end{table}

% \begin{figure}[H]
% \centering
% \includegraphics[width=0.95\linewidth]{figures/chaewon.jpg}
% \caption{Example of figure of my wife}
% \label{fig:my_wife}
% \end{figure}

% we can reference the figure \autoref{fig:my_wife}

\end{document}